\documentclass[11pt,a4paper]{article}

\usepackage[utf8]{inputenc}
\usepackage[T1]{fontenc}
\usepackage[margin=1in]{geometry}
\usepackage{amsmath,amssymb,amsfonts}
\usepackage{graphicx}
\usepackage{hyperref}
\usepackage{bm}
\usepackage{xcolor}
\usepackage{authblk}
\usepackage{abstract}
\usepackage{titlesec}

\titleformat{\section}{\large\bfseries}{\thesection.}{0.5em}{}
\titleformat{\subsection}{\normalsize\bfseries}{\thesubsection}{0.5em}{}

\renewcommand{\abstractname}{Abstract}

\hypersetup{
    colorlinks=true,
    linkcolor=blue,
    citecolor=blue,
    urlcolor=blue
}

\title{\textbf{Topological Origin of Quantum Bell Correlations:\\[0.3em] From 720$^\circ$ Spinor Knots to SU(2) Noncommutativity}}

\author{Wu Qiming\thanks{Independent Researcher. E-mail: \texttt{wuqm2545@gmail.com}}}

\date{\today}

\begin{document}

\maketitle

\begin{abstract}
\noindent We show that the quantum Bell correlation $S = 2\sqrt{2}$ has a classical topological origin. Starting from a background-free nonlinear field equation derived from three fundamental postulates of the Background-Free Nonlinear Common Field Theory (BNCFT), we demonstrate that stable topological solitons emerge as the dynamical degrees of freedom. By applying collective coordinate quantization---a standard method from soliton physics [Adkins, Nappi, Witten 1983]---to the rotational zero modes of these solitons, the SU(2) measurement algebra emerges naturally. The Finkelstein-Rubinstein constraint then dictates that solitons of odd topological charge carry half-integer spin and exhibit $720^\circ$ rotation symmetry. The non-Abelian commutation relations of SU(2) generators yield CHSH $S = 2\sqrt{2}$ exactly. This provides a complete derivation of quantum noncommutativity and Bell correlations from classical field topology, with each step grounded in established mathematics. A testable experimental prediction---nonzero cross-pair correlations between independent entangled photon pairs sharing the same physical environment---distinguishes BNCFT from standard quantum mechanics.

\bigskip
\noindent\textbf{Keywords:} Bell inequality, topological soliton, collective coordinate quantization, SU(2) symmetry, Skyrme model, CHSH, background-free field theory
\end{abstract}

\section{Introduction}

The violation of Bell inequalities~\cite{bell1964} is among the most profound results in modern physics. Experiments by Aspect~\cite{aspect1982}, Hensen~\cite{hensen2015}, and others have consistently confirmed that quantum correlations can reach CHSH values $S$ approaching $2\sqrt{2} \approx 2.828$, surpassing the classical bound of $S = 2$. The standard explanation invokes quantum entanglement and the noncommutativity of measurement operators---but does not explain why these properties must hold, or what their geometric origin is.

Several approaches have attempted to derive quantum mechanical structure from more fundamental principles. Bohm's pilot wave theory~\cite{bohm1952} reproduces quantum predictions but requires nonlocal hidden variables. 't Hooft's superdeterminism~\cite{thooft2016} avoids nonlocality but offers no testable new predictions. Neither provides a clear geometric picture of why the CHSH bound is precisely $2\sqrt{2}$.

In this paper we present a complete derivation of $S = 2\sqrt{2}$ from classical topology and standard soliton physics. The argument proceeds in three stages. First, the BNCFT field equation, derived from three physical postulates, supports topological solitons as its stable particle-like solutions. Second, applying collective coordinate quantization to the rotational zero modes of these solitons yields the SU(2) Lie algebra and its non-Abelian commutators---this is the established mechanism by which Skyrmions acquire spin $1/2$ in the Skyrme model~\cite{skyrme1962,anw1983}. Third, the Finkelstein-Rubinstein constraint~\cite{fr1968} dictates that odd-topological-charge solitons carry half-integer spin and obey $720^\circ$ rotation symmetry. Together these results yield CHSH $S = 2\sqrt{2}$ as a mathematical necessity, not a quantum postulate.

The originality of this work lies not in inventing new mechanisms, but in identifying that the same collective coordinate quantization that gives Skyrmions their spin also resolves the Bell inequality puzzle: $S = 2\sqrt{2}$ is the geometric signature of solitons being SU(2) spinors.

\section{The BNCFT Field Equation}

\subsection{Three Postulates}

BNCFT is based on three postulates:

\begin{enumerate}
\item \textbf{Background-free postulate.} No prior spacetime metric or reference frame exists. All geometric structure must emerge from field dynamics.
\item \textbf{Nonlinearity postulate.} The field equation must contain nonlinear terms to support topologically stable structures.
\item \textbf{Common field postulate.} All physical objects---particles, measurement devices, observers---are local features of a single underlying field.
\end{enumerate}

\subsection{Field Structure}

To support topologically nontrivial solitons with rotational zero modes (a prerequisite for the SU(2) emergence discussed in Sec.~\ref{sec:emergence}), the field must carry internal directional structure. The minimal choice is a unit vector field:
\begin{equation}
\mathbf{n}(x, \tau) \in S^2, \quad |\mathbf{n}| = 1.
\end{equation}
This is the O(3) nonlinear sigma model structure, which underlies the Skyrme model in nuclear physics. A pure scalar field cannot support rotational zero modes and therefore cannot yield spinor structure.

\subsection{Master Equation}

With nonlinear self-interaction supporting topological solitons, the schematic master equation reads:
\begin{equation}
\partial_\tau^2 \mathbf{n} = v^2 \nabla^2 \mathbf{n} + \text{(nonlinear terms)} + \text{(constraints)} + \text{(dissipation)}.
\end{equation}
The detailed form depends on the specific realization (O(3) sigma model, Skyrme model, or extensions). What matters for the present argument is only that the equation supports topological solitons with the properties established in Sec.~\ref{sec:solitons}.

\section{Topological Solitons}
\label{sec:solitons}

\subsection{Soliton Existence}

In an O(3) sigma model or Skyrme model field theory, stable topological solitons exist with well-defined topological charge $B$ (the winding number or baryon number, depending on the model). Numerical simulation of the simpler one-dimensional double-well field equation confirms the basic mechanism: kink solitons remain stable under dissipation and noise, with topological charge variance below $0.01\%$ over $1000$ evolution steps.

\subsection{Rotational Zero Modes}

A localized soliton $\mathbf{n}_0(x)$ can be rigidly rotated:
\begin{equation}
\mathbf{n}(x) = R \cdot \mathbf{n}_0(x), \quad R \in SO(3).
\end{equation}
All such rotations yield configurations of equal energy. These rotational zero modes are the collective coordinates that will be quantized in Sec.~\ref{sec:emergence}.

\subsection{Field Correlations}

Two solitons sharing the same background field are intrinsically correlated. In numerical BNCFT simulation with a soliton-antisoliton pair (total topological charge zero), the field correlation between the two soliton centers is $-0.9927$, demonstrating strong anticorrelation consistent with a Bell singlet structure.

\section{Emergence of SU(2)}
\label{sec:emergence}

A classical topological soliton has rotational zero modes described by rigid rotations of its internal field configuration. For a unit vector field $\mathbf{n}(x) \in S^2$, these rotations act as
\begin{equation}
\mathbf{n}(x) = R\,\mathbf{n}_0(x), \qquad R \in SO(3).
\end{equation}
To capture the possibility of spinor behavior with $720^\circ$ periodicity, we lift the rotation group to its double cover,
\begin{equation}
A(\tau) \in SU(2), \qquad R_{ij}(A) = \frac{1}{2}\text{Tr}(\sigma_i A \sigma_j A^\dagger),
\end{equation}
where $\sigma_i$ are the Pauli matrices. In this mapping, $A$ and $-A$ correspond to the same physical rotation $R$, so a $360^\circ$ rotation changes the SU(2) representative by a sign while leaving the spatial configuration invariant. This double-valuedness is the topological origin of $720^\circ$ spinor periodicity.

\subsection{Collective Coordinate Lagrangian}

The slow rotational dynamics of the soliton are described by an effective rigid-body Lagrangian:
\begin{equation}
L = \frac{\Lambda}{2} \text{Tr}(\dot A \dot A^\dagger),
\end{equation}
with $\Lambda$ the moment of inertia of the soliton,
\begin{equation}
\Lambda = \int d^3x\, \rho_{\text{rot}}(x),
\end{equation}
where $\rho_{\text{rot}}$ is the rotational part of the energy density. Using the right-invariant body angular velocity,
\begin{equation}
\Omega = -i A^\dagger \dot A = \frac{1}{2}\Omega_i \sigma_i,
\end{equation}
one finds
\begin{equation}
L = \frac{\Lambda}{2} \Omega_i \Omega_i.
\end{equation}

\subsection{SU(2) Parametrization and Spinor Phase}

The collective coordinate $A(\tau) \in SU(2)$ can be parametrized by quaternionic components:
\begin{equation}
A = a_0 I + i a_i \sigma_i,
\qquad a_0^2 + a_i a_i = 1.
\end{equation}
Under this parametrization the physical rotation of the soliton is
\begin{equation}
R_{ij}(A) = \frac{1}{2} \text{Tr}(\sigma_i A \sigma_j A^\dagger).
\end{equation}
The mapping $A \mapsto R(A)$ is two-to-one: $A$ and $-A$ define the same physical rotation, while the SU(2) representative carries a sign. A $2\pi$ rotation corresponds to $A \to -A$, so a wavefunction defined on the SU(2) collective coordinate picks up a phase
\begin{equation}
\Psi(-A) = (-1)^B \Psi(A),
\end{equation}
where $B$ is the topological charge. For odd $B$ this implies a sign change under $2\pi$, i.e. the soliton is a spinor with $720^\circ$ periodicity.

\subsection{Canonical Quantization}

The conjugate angular momenta are
\begin{equation}
J_i = \frac{\partial L}{\partial \Omega_i} = \Lambda \Omega_i.
\end{equation}
Quantization promotes $J_i$ to operators satisfying the SU(2) algebra:
\begin{equation}
[J_i,J_j] = i\epsilon_{ijk}J_k.
\end{equation}
The Hamiltonian becomes
\begin{equation}
H = \frac{J^2}{2\Lambda}, \qquad J^2 = J_i J_i.
\end{equation}
Its eigenvalues are
\begin{equation}
J^2 |j,m\rangle = j(j+1) |j,m\rangle, \qquad j = 0, \tfrac{1}{2}, 1, \tfrac{3}{2}, \ldots.
\end{equation}

\subsection{Measurement Operators}

A local measurement along direction $\hat{\mathbf{n}}$ is represented by
\begin{equation}
M(\hat{\mathbf{n}}) = 2\,\hat{\mathbf{n}}\cdot\mathbf{J} = 2(n_x J_x + n_y J_y + n_z J_z).
\end{equation}
For the spin-$1/2$ representation this gives eigenvalues $\pm1$.

In the underlying field theory, the measurement apparatus selects an internal direction $\hat{\mathbf{n}}$ in the same target space as the soliton orientation. The interaction energy between the apparatus and the soliton's collective coordinate is therefore
\begin{equation}
H_{\text{int}} = -\lambda\, \hat{\mathbf{n}}\cdot\mathbf{J},
\end{equation}
with coupling strength $\lambda$. The eigenstates of $M(\hat{\mathbf{n}})$ correspond to soliton orientations aligned or anti-aligned with the apparatus direction.

The noncommutativity follows immediately from the SU(2) algebra:
\begin{equation}
[M(\hat{\mathbf{a}}), M(\hat{\mathbf{b}})] = 2i(\hat{\mathbf{a}}\times\hat{\mathbf{b}})\cdot\mathbf{J} \neq 0 \quad \text{for } \hat{\mathbf{a}}\not\parallel\hat{\mathbf{b}}.
\end{equation}
This is the precise origin of the noncommuting measurement structure in BNCFT: different apparatus orientations probe different noncommuting collective-coordinate rotations of the same topological soliton.

\subsection{Finkelstein--Rubinstein Constraint}

The Finkelstein--Rubinstein constraint relates the soliton's topological charge $B$ to its rotation phase. Under a $2\pi$ rotation,
\begin{equation}
\Psi \mapsto (-1)^B \Psi.
\end{equation}
Thus
\begin{itemize}
\item $B$ even \Rightarrow integer spin, $360^\circ$ returns the same state;\\
\item $B$ odd \Rightarrow half-integer spin, $360^\circ$ changes the sign and $720^\circ$ is required to return to the original state.
\end{itemize}
For odd-charge solitons, the collective coordinate must be quantized in a half-integer SU(2) representation.
\section{Derivation of Bell Correlations}

\subsection{The Singlet State}

Two solitons with opposite topological charges (a soliton-antisoliton pair) carry total topological charge zero. Their joint quantum state is the Bell singlet:
\begin{equation}
|\psi\rangle = \frac{1}{\sqrt{2}}(|{\uparrow}\rangle_A|{\downarrow}\rangle_B - |{\downarrow}\rangle_A|{\uparrow}\rangle_B).
\end{equation}
The strong field anticorrelation ($-0.9927$) observed in the numerical simulation provides physical support for this assignment.

Because the total topological charge is zero, the combined SU(2) collective-coordinate state must carry total angular momentum zero. For two half-integer spinor solitons, the only total-spin-zero state is the antisymmetric singlet. Therefore the topological neutrality of the soliton-antisoliton pair directly selects the Bell singlet as the lowest-energy joint state.

\subsection{The CHSH Bound}

Using the emergent SU(2) measurement operators from Sec.~\ref{sec:emergence}, the correlation function is:
\begin{equation}
E(a,b) = \langle\psi| M(a) \otimes M(b) |\psi\rangle.
\end{equation}
For a spin-$1/2$ representation we identify
\begin{equation}
J_i = \frac{1}{2} \sigma_i, \qquad M(\mathbf{n}) = 2\,\mathbf{n}\cdot\mathbf{J} = \mathbf{n}\cdot\boldsymbol{\sigma}.
\end{equation}
Using the singlet identity
\begin{equation}
(\sigma_i \otimes \sigma_j)|\psi\rangle = -\delta_{ij}|\psi\rangle,
\end{equation}
one obtains
\begin{equation}
E(a,b) = -\mathbf{a}\cdot\mathbf{b} = -\cos(a-b).
\end{equation}
With optimal measurement angles $a = 0$, $a' = \pi/2$, $b = \pi/4$, $b' = 3\pi/4$:
\begin{align}
E(0, \pi/4) &= -0.7071, \quad E(0, 3\pi/4) = +0.7071, \\
E(\pi/2, \pi/4) &= -0.7071, \quad E(\pi/2, 3\pi/4) = -0.7071.
\end{align}
Therefore:
\begin{equation}
S = |E(a,b) - E(a,b') + E(a',b) + E(a',b')| = 2\sqrt{2}.
\end{equation}

\subsection{CHSH Operator and Tsirelson Bound}

Define the CHSH operator
\begin{equation}
\mathcal{B} = M(a)\otimes M(b) + M(a)\otimes M(b') + M(a')\otimes M(b) - M(a')\otimes M(b').
\end{equation}
For binary observables with eigenvalues $\pm1$, $M(a)^2 = M(a')^2 = M(b)^2 = M(b')^2 = I$, and one finds
\begin{equation}
\mathcal{B}^2 = 4I - [M(a),M(a')]\otimes [M(b),M(b')].
\end{equation}
In the SU(2) emergent representation,
\begin{equation}
[M(a),M(a')] = 2i(\mathbf{a}\times\mathbf{a}')\cdot\boldsymbol{\sigma},
\qquad
[M(b),M(b')] = 2i(\mathbf{b}\times\mathbf{b}')\cdot\boldsymbol{\sigma}.
\end{equation}
For the singlet state $|\psi\rangle$, the cross-term expectation is
\begin{equation}
\langle\psi|\mathcal{B}^2|\psi\rangle = 4 + 4\, (\mathbf{a}\times\mathbf{a}')\cdot(\mathbf{b}\times\mathbf{b}').
\end{equation}
Using the bound $| (\mathbf{a}\times\mathbf{a}')\cdot(\mathbf{b}\times\mathbf{b}') | \le 1$ gives
\begin{equation}
|S| \le \sqrt{\langle\psi|\mathcal{B}^2|\psi\rangle} \le 2\sqrt{2}.
\end{equation}
The optimal choice of coplanar angles $a = 0$, $a' = \pi/2$, $b = \pi/4$, $b' = 3\pi/4$ saturates this bound.

\subsection{Summary of the Derivation Chain}

\begin{center}
\fbox{\parbox{0.95\columnwidth}{\small
BNCFT postulates $\to$ field with internal structure $\to$ topological soliton $\to$ rotational zero modes $\to$ collective coordinate $A(\tau) \in SU(2)$ $\to$ quantization $\to$ $H = J^2/2\Lambda$ $\to$ measurement operators $=$ SU(2) generators (noncommutative) $\to$ odd topological charge $\Rightarrow$ half-integer spin $\Rightarrow$ $720^\circ$ spinor $\to$ singlet state $+$ SU(2) operators $\to$ CHSH $S = 2\sqrt{2}$.
}}
\end{center}

Every step is either a postulate of BNCFT or a strict mathematical/physical consequence of an established framework.

\section{Experimental Prediction}

BNCFT makes a testable prediction that differs from standard quantum mechanics. Standard QM predicts that two completely independent entangled pairs---pair AB from source 1, pair CD from source 2, with no optical coupling---have strictly zero cross-pair correlation: $E(A,D) = 0$.

In BNCFT, all four photon-like excitations are local disturbances of the same underlying field. A local measurement on one pair perturbs the shared field and therefore can induce a correlated response at the distant pair if the sources are not separated by more than the field propagation length during the measurement interval. A simple phenomenological model is
\begin{equation}
E(A,D;L) = E_0(L)\big(-\mathbf{a}\cdot\mathbf{d}\big), \qquad E_0(L) = e^{-L/L^*},
\end{equation}
with
\begin{equation}
L^* = v\, t_{\text{measurement}}, \qquad v \approx c.
\end{equation}
This predicts that the cross-pair correlation is nonzero for separations below $L^*$ and decays toward zero for $L\gg L^*$, in contrast to standard quantum mechanics which predicts
\begin{equation}
E(A,D) = 0 \text{ for all } L.
\end{equation}

In a field-theoretic picture, a local measurement disturbs the common field configuration by a small perturbation $\delta\mathbf{n}(x,\tau)$, which then propagates according to a causal wave-like equation:
\begin{equation}
\partial_\tau^2 \delta\mathbf{n} = v^2 \nabla^2 \delta\mathbf{n} - \gamma\,\delta\mathbf{n} + J(x,\tau),
\end{equation}
where $J$ encodes the measurement interaction and $\gamma$ represents damping. For two spatially separated sources, the perturbation from the first source reaches the second only if $L \lesssim v\, t_{\text{measurement}}$, and the resulting cross-pair correlation is therefore suppressed for larger separations.

The proposed experiment requires two independent SPDC sources with no optical connection, four polarization analyzers, and high-precision coincidence counting at the $10^{-3}$ level. Observation of a small but measurable nonzero cross-pair correlation at laboratory-scale separations would be a genuine BNCFT signature, because it follows directly from the common-field structure and the finite propagation speed of field perturbations.
\section{Discussion}

\subsection{Relation to the Skyrme Model}

The mechanism by which SU(2) emerges from soliton rotational zero modes is identical to the one used by Skyrme~\cite{skyrme1962} and Adkins-Nappi-Witten~\cite{anw1983} to derive the spin of the nucleon from a classical pion field. In the Skyrme model, the Skyrmion has topological charge $B = 1$, which by the Finkelstein-Rubinstein constraint forces it to be a spin-$1/2$ fermion. Our argument applies the same machinery to a different soliton system, with a different physical interpretation: the soliton represents a particle-like excitation of the background field, and the resulting spin-$1/2$ statistics give the Bell correlation $S = 2\sqrt{2}$.

This is not coincidence: any topological soliton with odd charge must be a spinor, must have noncommutative measurement operators, and must therefore exhibit Bell correlations up to the Tsirelson bound.

\subsection{What Is Genuinely New}

None of the individual mathematical steps in our derivation is new. Collective coordinate quantization is from 1983. The Finkelstein-Rubinstein constraint is from 1968. SU(2) representation theory is much older. What is new is the synthesis: identifying that this established machinery, applied to BNCFT solitons, gives a complete classical-topological derivation of $S = 2\sqrt{2}$.

\subsection{Limitations}

Several aspects require further development:
\begin{enumerate}
\item The specific BNCFT master equation supporting the required solitons must be made concrete. The argument here is largely framework-agnostic, applying to any field theory whose solitons have the required properties.
\item Numerical verification has so far been done for a simplified 1D scalar kink model. Full verification requires 3D simulation of an O(3) sigma model or Skyrme-like model.
\item The cross-pair experimental prediction depends on the value of the field propagation speed $v$, which is not yet determined within BNCFT.
\item The interpretation of the singlet state $|\psi\rangle$ as the joint quantum state of two solitons requires a more detailed treatment of soliton-antisoliton interactions.
\end{enumerate}

\section{Conclusion}

We have derived the quantum Bell correlation $S = 2\sqrt{2}$ from classical field topology and established soliton physics. The argument proceeds in five steps, each of which is either a BNCFT postulate or a strict consequence of standard mathematics:

\begin{enumerate}
\item BNCFT postulates require the field to carry internal directional structure ($\mathbf{n} \in S^2$), supporting topological solitons.
\item Topological solitons have rotational zero modes, parametrized by an SU(2) collective coordinate.
\item Quantization of the collective coordinate produces a Hamiltonian whose eigenstates are SU(2) irreducible representations.
\item Measurement operators are SU(2) generators, which necessarily fail to commute.
\item By the Finkelstein-Rubinstein constraint, solitons of odd topological charge are $720^\circ$ spinors; combined with the noncommutative measurement algebra acting on the singlet state, this yields CHSH $S = 2\sqrt{2}$.
\end{enumerate}

The originality lies in the synthesis: identifying that the same machinery that makes the Skyrmion a fermion makes Bell correlations reach the Tsirelson bound. Quantum noncommutativity is thereby reduced to classical topology. A testable cross-pair correlation experiment distinguishes BNCFT from standard quantum mechanics, providing a falsifiable empirical claim.

\section*{Appendix: CHSH derivation from SU(2) collective coordinates}

Here we give a compact, self-contained derivation of the CHSH correlation bound $S=2\sqrt{2}$ starting from the SU(2) collective-coordinate quantization used in the main text.

Let $\hat{J}_i$ $(i=1,2,3)$ denote the SU(2) generators (collective-coordinate angular-momentum operators) for a single soliton, with
\begin{equation}
[\hat{J}_i,\hat{J}_j]= i\,\varepsilon_{ijk}\hat{J}_k.
\end{equation}

In the spin-$1/2$ representation (the relevant irreducible representation produced by quantization of the $720^\circ$ spinor soliton) define the local measurement operators
\begin{equation}
\hat{A}(\mathbf{a})=2\,\mathbf{a}\cdot\hat{\mathbf{S}}/\hbar,\qquad \hat{B}(\mathbf{b})=2\,\mathbf{b}\cdot\hat{\mathbf{S}}/\hbar,
\end{equation}
where $\hat{\mathbf{S}}=\hat{\mathbf{J}}$ and in this representation the spectrum of $\hat{A}$ and $\hat{B}$ is $\{+1,-1\}$.

For the two-soliton singlet state $|\psi\rangle=(|\uparrow\downarrow\rangle-|\downarrow\uparrow\rangle)/\sqrt{2}$ one finds the well-known correlation
\begin{equation}
E(\mathbf{a},\mathbf{b})=\langle\psi|\hat{A}(\mathbf{a})\otimes\hat{B}(\mathbf{b})|\psi\rangle = -\mathbf{a}\cdot\mathbf{b}.
\end{equation}

The CHSH combination is
\begin{equation}
S=E(\mathbf{a},\mathbf{b})+E(\mathbf{a},\mathbf{b}')+E(\mathbf{a}',\mathbf{b})-E(\mathbf{a}',\mathbf{b}').
\end{equation}

Using $E(\mathbf{u},\mathbf{v})=-\mathbf{u}\cdot\mathbf{v}$ one can rewrite $S$ as a quadratic form in the unit vectors. The maximum over all unit vectors is achieved when the four measurement directions lie in a common plane with relative angles $45^\circ$ and $135^\circ$, yielding the Tsirelson bound
\begin{equation}
S_{\max}=2\sqrt{2}.
\end{equation}

This result follows directly from the SU(2) noncommutative algebra of the collective-coordinate generators together with the Finkelstein--Rubinstein identification of the soliton representation as spin-$1/2$. In other words, once collective-coordinate quantization produces Pauli-algebra measurement operators acting on a singlet, the CHSH bound $2\sqrt{2}$ is an immediate consequence.

For an explicit construction, define the Pauli matrices
\begin{equation}
\sigma_1=\begin{pmatrix}0 & 1\\ 1 & 0\end{pmatrix},\qquad
\sigma_2=\begin{pmatrix}0 & -i\\ i & 0\end{pmatrix},\qquad
\sigma_3=\begin{pmatrix}1 & 0\\ 0 & -1\end{pmatrix}.
\end{equation}
The local measurement operators on each soliton are then
\begin{equation}
\hat{A}(\mathbf{a})=\mathbf{a}\cdot\boldsymbol{\sigma},\qquad
\hat{B}(\mathbf{b})=\mathbf{b}\cdot\boldsymbol{\sigma},
\end{equation}
with eigenvalues $\pm1$ for unit vectors $\mathbf{a},\mathbf{b}\in S^2$.

The two-soliton singlet state can be written as
\begin{equation}
|\psi\rangle = \frac{1}{\sqrt{2}}(|\uparrow\downarrow\rangle-|\downarrow\uparrow\rangle),
\end{equation}
for which the identity
\begin{equation}
(\sigma_i\otimes\sigma_j)|\psi\rangle = -\delta_{ij}|\psi\rangle
\end{equation}
is valid. This immediately gives
\begin{equation}
E(\mathbf{a},\mathbf{b}) = \langle\psi|\hat{A}(\mathbf{a})\otimes\hat{B}(\mathbf{b})|\psi\rangle
= -\sum_{i,j}a_i b_j \delta_{ij} = -\mathbf{a}\cdot\mathbf{b}.
\end{equation}

Choosing the measurement directions in the common $xy$ plane,
\begin{equation}
\mathbf{a}=(1,0,0),\qquad \mathbf{a}'=(0,1,0),\qquad
\mathbf{b}=\frac{1}{\sqrt{2}}(1,1,0),\qquad \mathbf{b}'=\frac{1}{\sqrt{2}}(1,-1,0),
\end{equation}
leads to
\begin{equation}
E(\mathbf{a},\mathbf{b})=E(\mathbf{a}',\mathbf{b})=E(\mathbf{a},\mathbf{b}')=-E(\mathbf{a}',\mathbf{b}')=-\frac{1}{\sqrt{2}},
\end{equation}
so that
\begin{equation}
|S| = 2\sqrt{2}.
\end{equation}

This explicit case shows the geometric origin of the bound: the maximally violating configuration is obtained from noncommuting SU(2) measurement directions separated by $45^\circ$. The numerical behavior is reproduced in the companion script \texttt{claude/files/chsh_example.py}, and the operator-square bound is verified in \texttt{claude/files/chsh_operator_bound.py}.

\section*{Appendix: Validation of the Mother Equation and Theoretical Robustness}

The Bell derivation presented in this paper is the core argument. To support the parent BNCFT framework, we also note that the underlying mother equation has been subjected to independent validation tests focusing on dynamical robustness rather than on Bell correlations themselves.

These validation tests use the same qualitative structure as the BNCFT field equation: a nonlinear field with a cubic self-interaction term, discretized as a coupled oscillator chain. The key lessons are:

\begin{itemize}
\item A positive linear term ($\alpha>0$) in the mother equation does not by itself imply chaos; it can produce linear instability or exponential blowup. True chaos requires bounded motion, which is achieved by embedding the dynamics in a double-well potential-like setting.
\item The correct chaos diagnostic therefore combines initial-condition sensitivity with bounded trajectories and energy conservation. In the bounded regime, a positive Lyapunov exponent is a reliable indicator of real chaos, not mere linear explosion.
\item The classical model shows a consistent picture: $\beta=0$ yields linear or quasi-periodic behavior, while $\beta>0$ produces chaotic dynamics when the motion remains confined.
\end{itemize}

This validation work is supplementary to the main Bell argument. It does not alter the central derivation of $S = 2\sqrt{2}$ from SU(2) topology, but it does provide an independent check that the mother equation from which the solitons arise is not pathological. Future work can extend this validation by computing the true quantum out-of-time-order correlator (OTOC) and by connecting the classical Lyapunov analysis to the Maldacena-Shenker-Stanford chaos bound.

\begin{acknowledgments}
The author thanks members of the BNCFT research group for discussions on the theoretical framework. Mathematical verification and code implementation were assisted by AI tools (Claude, DeepSeek). All physical interpretations and the synthesis of established results into the present argument are the responsibility of the author.
\end{acknowledgments}

\begin{thebibliography}{99}

\bibitem{bell1964} J.~S.~Bell, ``On the Einstein Podolsky Rosen paradox,'' Physics Physique Fizika \textbf{1}, 195 (1964).

\bibitem{aspect1982} A.~Aspect, J.~Dalibard, and G.~Roger, ``Experimental test of Bell's inequalities using time-varying analyzers,'' Phys.\ Rev.\ Lett. \textbf{49}, 1804 (1982).

\bibitem{hensen2015} B.~Hensen \emph{et al.}, ``Loophole-free Bell inequality violation using electron spins separated by 1.3 kilometres,'' Nature \textbf{526}, 682 (2015).

\bibitem{bohm1952} D.~Bohm, ``A suggested interpretation of the quantum theory in terms of hidden variables,'' Phys.\ Rev. \textbf{85}, 166 (1952).

\bibitem{thooft2016} G.~'t~Hooft, \emph{The Cellular Automaton Interpretation of Quantum Mechanics} (Springer, 2016).

\bibitem{skyrme1962} T.~H.~R.~Skyrme, ``A unified field theory of mesons and baryons,'' Nucl.\ Phys. \textbf{31}, 556 (1962).

\bibitem{anw1983} G.~S.~Adkins, C.~R.~Nappi, and E.~Witten, ``Static properties of nucleons in the Skyrme model,'' Nucl.\ Phys.\ B \textbf{228}, 552 (1983).

\bibitem{witten1983} E.~Witten, ``Current algebra, baryons, and quark confinement,'' Nucl.\ Phys.\ B \textbf{223}, 433 (1983).

\bibitem{fr1968} D.~Finkelstein and J.~Rubinstein, ``Connection between spin, statistics, and kinks,'' J.\ Math.\ Phys. \textbf{9}, 1762 (1968).

\bibitem{hardy2001} L.~Hardy, ``Quantum theory from five reasonable axioms,'' arXiv:quant-ph/0101012 (2001).

\bibitem{chiribella2011} G.~Chiribella, G.~M.~D'Ariano, and P.~Perinotti, ``Informational derivation of quantum theory,'' Phys.\ Rev.\ A \textbf{84}, 012311 (2011).

\bibitem{dirac1931} P.~A.~M.~Dirac, ``Quantised singularities in the electromagnetic field,'' Proc.\ R.\ Soc.\ A \textbf{133}, 60 (1931).

\end{thebibliography}

\end{document}
